Parkinson's disease (PD) emerges from a prolonged prodromal phase characterized by non-motor features—REM sleep behavior disorder (RBD), hyposmia, constipation—preceding motor symptom onset by years to decades\cite{postuma2015prodromal,berg2015time}. Identifying prodromal individuals at high risk for phenoconversion to clinical PD is critical for: (1) early intervention before substantial neurodegeneration; (2) enriching preventive trials with fast-converting participants; (3) personalized clinical counseling on conversion probability and timing\cite{marek2013mdsuppcc,postuma2019risk}.

Prodromal cohort studies have identified risk factors for conversion: RBD confers 3-fold increased risk\cite{postuma2015prodromal}, hyposmia 2-3 fold\cite{ponsen2004idiopathic}, abnormal DaTscan 3-4 fold\cite{jennings2017conversion}, and GBA mutations 5-10 fold\cite{alcalay2012cognitive,iwaki2019genetic}. However, most studies report univariate associations or simple risk scores\cite{berg2015time}, lacking personalized time-to-conversion predictions. Cox proportional hazards models\cite{cox1972regression}, while standard for survival analysis, assume linear hazard relationships and cannot capture complex biomarker interactions or time-varying effects.

Deep learning offers a solution. DeepSurv\cite{katzman2018deepsurv}, a deep neural network for survival analysis, learns nonlinear risk functions from high-dimensional covariates, outperforming Cox models in cancer\cite{katzman2018deepsurv}, cardiovascular disease\cite{nagpal2021deep}, and Alzheimer's disease\cite{nguyen2021deep}. However, DeepSurv has not been applied to prodromal PD conversion, where complex interactions between genetics, biomarkers, and clinical features likely govern conversion risk.

We developed and validated \deepsurv{} models for personalized prodromal-to-clinical PD conversion prediction in the Parkinson's Progression Markers Initiative (\ppmi). Our objectives were: (1) Compare \deepsurv{} vs. Cox models for discrimination and calibration; (2) Identify key baseline predictors and their hazard relationships; (3) Define biomarker-based risk strata for clinical utility; (4) Simulate preventive trial enrichment to quantify sample size reductions. We hypothesized that \deepsurv's nonlinear modeling would improve prediction over Cox, and that high-risk strata would enable efficient trial designs.
